%%%%%%%%%%%%%%%%%%%%%%%%%%%%%%%%%%%%%%%%%%%%%%%%%%%%%%%%%%%%
% 0.10 MATHEMATICAL WRITING AND PROBLEM SOLVING
% The Composition of Advanced Mathematics
%%%%%%%%%%%%%%%%%%%%%%%%%%%%%%%%%%%%%%%%%%%%%%%%%%%%%%%%%%%%

\section{Mathematical Writing and Problem Solving}
\label{sec:mathematical-writing-problem-solving}

\prerequisite{
Familiarity with mathematical notation, sets, functions, logic, and proof
methods from the preceding sections of Chapter 0.
}

\learningobjectives{
By the end of this section, the reader should be able to:
\begin{itemize}
    \item read mathematical writing actively rather than mechanically;
    \item use definitions as working mathematical tools;
    \item write mathematical expressions as parts of complete sentences;
    \item introduce variables and notation clearly;
    \item organize displayed equations and derivations effectively;
    \item translate between mathematical notation and spoken language;
    \item identify the givens, unknowns, conditions, and goals of a problem;
    \item apply a systematic problem-solving process;
    \item use diagrams, tables, patterns, and special cases strategically;
    \item work both forward from known information and backward from a goal;
    \item use estimation, units, domains, and special cases to check results;
    \item distinguish exploratory work from a polished mathematical solution;
    \item write clear proofs and worked solutions;
    \item recognize common weaknesses in mathematical communication.
\end{itemize}
}

%%%%%%%%%%%%%%%%%%%%%%%%%%%%%%%%%%%%%%%%%%%%%%%%%%%%%%%%%%%%
\subsection{Mathematics as Communication}
\label{subsec:mathematics-as-communication}
%%%%%%%%%%%%%%%%%%%%%%%%%%%%%%%%%%%%%%%%%%%%%%%%%%%%%%%%%%%%

Mathematics is not merely a collection of calculations.

It is a language for expressing definitions, relationships, patterns,
arguments, structures, and conclusions precisely.

A correct calculation can still be difficult to understand if its notation is
unclear or its reasoning is poorly organized.

Likewise, a well-written mathematical argument allows another reader to
reconstruct not only what was done, but why each step is valid.

\begin{conceptbox}
Good mathematical writing aims for three qualities:
\begin{itemize}
    \item \textbf{correctness} --- every statement is mathematically valid;
    \item \textbf{clarity} --- the reader can understand what each symbol and step means;
    \item \textbf{precision} --- assumptions, domains, definitions, and conclusions are stated accurately.
\end{itemize}
\end{conceptbox}

%%%%%%%%%%%%%%%%%%%%%%%%%%%%%%%%%%%%%%%%%%%%%%%%%%%%%%%%%%%%
\subsection{Reading Mathematics Actively}
\label{subsec:reading-mathematics-actively}
%%%%%%%%%%%%%%%%%%%%%%%%%%%%%%%%%%%%%%%%%%%%%%%%%%%%%%%%%%%%

Mathematical reading is usually slower than ordinary reading.

A single displayed equation may contain several relationships that must be
interpreted before continuing.

When reading mathematics, ask:

\begin{itemize}
    \item What objects are being discussed?
    \item What do the symbols represent?
    \item What is assumed?
    \item What is being claimed?
    \item Which definitions are being used?
    \item Why does one step follow from the previous step?
    \item Is the statement an example, definition, theorem, or conclusion?
\end{itemize}

Reading actively often requires stopping and reconstructing omitted
intermediate reasoning.

%%%%%%%%%%%%%%%%%%%%%%%%%%%%%%%%%%%%%%%%%%%%%%%%%%%%%%%%%%%%
\subsection{Reading Symbols as Meaning}
\label{subsec:reading-symbols-as-meaning}
%%%%%%%%%%%%%%%%%%%%%%%%%%%%%%%%%%%%%%%%%%%%%%%%%%%%%%%%%%%%

Symbols should not be read merely as shapes.

For example,

\[ x\in A \]

should be understood as

\begin{center}
``\(x\) is an element of \(A\).''
\end{center}

Similarly,

\[ A\subseteq B \]

means

\begin{center}
``every element of \(A\) is also an element of \(B\).''
\end{center}

The symbolic form is compact, but the underlying meaning should remain
accessible.

%%%%%%%%%%%%%%%%%%%%%%%%%%%%%%%%%%%%%%%%%%%%%%%%%%%%%%%%%%%%
\subsection{Reading Mathematical Structure}
\label{subsec:reading-mathematical-structure}
%%%%%%%%%%%%%%%%%%%%%%%%%%%%%%%%%%%%%%%%%%%%%%%%%%%%%%%%%%%%

Consider

\[ f(x)=\frac{x^2-1}{x-1}. \]

A reader should notice more than the numerator and denominator.

The expression has a restriction:

\[ x\neq1. \]

The numerator factors:

\[ x^2-1=(x-1)(x+1). \]

Thus, for \(x\neq1\),

\[ f(x)=x+1. \]

The restriction remains important even after simplification.

Active reading identifies both algebraic structure and domain information.

%%%%%%%%%%%%%%%%%%%%%%%%%%%%%%%%%%%%%%%%%%%%%%%%%%%%%%%%%%%%
\subsection{Definitions Are Working Tools}
\label{subsec:definitions-working-tools}
%%%%%%%%%%%%%%%%%%%%%%%%%%%%%%%%%%%%%%%%%%%%%%%%%%%%%%%%%%%%

Definitions should not be treated as vocabulary to memorize and then ignore.

They tell us exactly what must be shown.

Suppose one wants to prove that a function is injective.

The definition says that one must establish

\[ f(x_1)=f(x_2)\Longrightarrow x_1=x_2. \]

That definition immediately suggests the structure of the proof.

Similarly, to prove

\[ A\subseteq B, \]

the definition tells us to begin with an arbitrary

\[ x\in A \]

and establish

\[ x\in B. \]

%%%%%%%%%%%%%%%%%%%%%%%%%%%%%%%%%%%%%%%%%%%%%%%%%%%%%%%%%%%%
\subsection{Unpacking Definitions}
\label{subsec:unpacking-definitions}
%%%%%%%%%%%%%%%%%%%%%%%%%%%%%%%%%%%%%%%%%%%%%%%%%%%%%%%%%%%%

A useful mathematical habit is to replace terminology with its precise
definition.

For example:

\[
\begin{array}{c|l}
\text{Term} & \text{Useful Mathematical Form}\\ \hline
n\text{ is even} & n=2k\text{ for some }k\in\Z\\
n\text{ is odd} & n=2k+1\text{ for some }k\in\Z\\
A\subseteq B & x\in A\Rightarrow x\in B\\
f\text{ injective} & f(x_1)=f(x_2)\Rightarrow x_1=x_2\\
f\text{ surjective} & \forall y\in B,\;\exists x\in A,\;f(x)=y
\end{array}
\]

Definitions convert abstract words into usable mathematical conditions.

%%%%%%%%%%%%%%%%%%%%%%%%%%%%%%%%%%%%%%%%%%%%%%%%%%%%%%%%%%%%
\subsection{Mathematical Sentences}
\label{subsec:mathematical-sentences}
%%%%%%%%%%%%%%%%%%%%%%%%%%%%%%%%%%%%%%%%%%%%%%%%%%%%%%%%%%%%

Mathematical expressions should usually function as parts of grammatical
sentences.

For example:

\begin{center}
Since \(x>0\), we have \(x^2>0\).
\end{center}

This is preferable to an unexplained sequence such as

\[ x>0,\qquad x^2>0. \]

The second version contains mathematical information but does not communicate
the logical relationship between the statements.

%%%%%%%%%%%%%%%%%%%%%%%%%%%%%%%%%%%%%%%%%%%%%%%%%%%%%%%%%%%%
\subsection{Symbols Do Not Replace Explanation}
\label{subsec:symbols-not-explanation}
%%%%%%%%%%%%%%%%%%%%%%%%%%%%%%%%%%%%%%%%%%%%%%%%%%%%%%%%%%%%

Symbolic notation is powerful because it compresses information.

However, excessive symbolic compression can obscure reasoning.

For example,

\[ x\in A\Rightarrow x\in B\Rightarrow A\subseteq B \]

may be understandable in context, but a proof is clearer when the role of the
arbitrary element is stated:

\begin{center}
Let \(x\in A\) be arbitrary. If \(x\in A\), then \(x\in B\). Therefore every
element of \(A\) belongs to \(B\), so \(A\subseteq B\).
\end{center}

%%%%%%%%%%%%%%%%%%%%%%%%%%%%%%%%%%%%%%%%%%%%%%%%%%%%%%%%%%%%
\subsection{Introducing Variables}
\label{subsec:introducing-variables}
%%%%%%%%%%%%%%%%%%%%%%%%%%%%%%%%%%%%%%%%%%%%%%%%%%%%%%%%%%%%

A variable should be introduced before it is used.

Instead of writing

\[ n=2k \]

without context, write:

\begin{center}
Let \(n\) be an even integer. Then \(n=2k\) for some \(k\in\Z\).
\end{center}

The reader now knows:

\begin{itemize}
    \item what \(n\) represents;
    \item why \(k\) exists;
    \item what set \(k\) belongs to.
\end{itemize}

%%%%%%%%%%%%%%%%%%%%%%%%%%%%%%%%%%%%%%%%%%%%%%%%%%%%%%%%%%%%
\subsection{Declare the Domain}
\label{subsec:declare-domain}
%%%%%%%%%%%%%%%%%%%%%%%%%%%%%%%%%%%%%%%%%%%%%%%%%%%%%%%%%%%%

A variable's domain can fundamentally change the meaning of a statement.

Consider

\[ x^2=2. \]

Over \(\R\), this equation has solutions

\[ x=\pm\sqrt2. \]

Over \(\Q\), it has no solutions.

Therefore statements such as

\begin{center}
``Let \(x\) be a real number''
\end{center}

or

\begin{center}
``Let \(n\in\Z\)''
\end{center}

often carry essential mathematical information.

%%%%%%%%%%%%%%%%%%%%%%%%%%%%%%%%%%%%%%%%%%%%%%%%%%%%%%%%%%%%
\subsection{Introduce Specialized Notation}
\label{subsec:introduce-specialized-notation}
%%%%%%%%%%%%%%%%%%%%%%%%%%%%%%%%%%%%%%%%%%%%%%%%%%%%%%%%%%%%

When notation is created for a local argument, explain it.

For example:

\begin{center}
Let \(E\) denote the set of even integers.
\end{center}

Then one may write

\[ E=\{2k:k\in\Z\}. \]

The notation should help the reader rather than force the reader to guess its
meaning.

%%%%%%%%%%%%%%%%%%%%%%%%%%%%%%%%%%%%%%%%%%%%%%%%%%%%%%%%%%%%
\subsection{Use Notation Consistently}
\label{subsec:notation-consistency}
%%%%%%%%%%%%%%%%%%%%%%%%%%%%%%%%%%%%%%%%%%%%%%%%%%%%%%%%%%%%

Once a symbol has been assigned a meaning, that meaning should remain stable
throughout the argument unless explicitly redefined.

If \(n\) denotes an integer at the beginning of a proof, it should not later
become a real-valued function without explanation.

Consistency reduces unnecessary cognitive effort.

%%%%%%%%%%%%%%%%%%%%%%%%%%%%%%%%%%%%%%%%%%%%%%%%%%%%%%%%%%%%
\subsection{Displayed Mathematics}
\label{subsec:displayed-mathematics}
%%%%%%%%%%%%%%%%%%%%%%%%%%%%%%%%%%%%%%%%%%%%%%%%%%%%%%%%%%%%

Important equations are often placed on their own line:

\[ ax+b=0. \]

Displayed mathematics is useful when:

\begin{itemize}
    \item an equation is central to the discussion;
    \item several transformations must be shown;
    \item the expression is too long for comfortable inline reading;
    \item the reader needs to refer visually to the equation.
\end{itemize}

Short expressions can remain inline when displaying them would interrupt the
flow unnecessarily.

%%%%%%%%%%%%%%%%%%%%%%%%%%%%%%%%%%%%%%%%%%%%%%%%%%%%%%%%%%%%
\subsection{Aligned Derivations}
\label{subsec:aligned-derivations}
%%%%%%%%%%%%%%%%%%%%%%%%%%%%%%%%%%%%%%%%%%%%%%%%%%%%%%%%%%%%

When several equations form one calculation, alignment can reveal the
structure.

For example,

\[
\begin{aligned}
(x+3)^2
&=(x+3)(x+3)\\
&=x^2+3x+3x+9\\
&=x^2+6x+9.
\end{aligned}
\]

The equality signs align vertically, making the chain of equivalent
expressions easy to follow.

%%%%%%%%%%%%%%%%%%%%%%%%%%%%%%%%%%%%%%%%%%%%%%%%%%%%%%%%%%%%
\subsection{Justifying Important Steps}
\label{subsec:justifying-steps}
%%%%%%%%%%%%%%%%%%%%%%%%%%%%%%%%%%%%%%%%%%%%%%%%%%%%%%%%%%%%

Not every elementary algebraic step needs a paragraph of explanation.

However, important transitions should be justified.

For example,

\[ a^2=b^2 \]

does not imply only

\[ a=b. \]

The correct conclusion over the real numbers is

\[ a=b\qquad\text{or}\qquad a=-b. \]

A written solution should make such logical branches visible.

%%%%%%%%%%%%%%%%%%%%%%%%%%%%%%%%%%%%%%%%%%%%%%%%%%%%%%%%%%%%
\subsection{Equality Signs Must Mean Equality}
\label{subsec:equality-sign-writing}
%%%%%%%%%%%%%%%%%%%%%%%%%%%%%%%%%%%%%%%%%%%%%%%%%%%%%%%%%%%%

The equality symbol should connect expressions that are actually equal.

A common poor habit is to write

\[ x^2-4=0=x^2=4=x=2. \]

This is not a valid chain of equalities.

A clearer derivation is

\[
\begin{aligned}
x^2-4&=0\\
x^2&=4\\
x&=\pm2.
\end{aligned}
\]

Each line now follows meaningfully from the previous one.

%%%%%%%%%%%%%%%%%%%%%%%%%%%%%%%%%%%%%%%%%%%%%%%%%%%%%%%%%%%%
\subsection{The Problem-Solving Process}
\label{subsec:problem-solving-process}
%%%%%%%%%%%%%%%%%%%%%%%%%%%%%%%%%%%%%%%%%%%%%%%%%%%%%%%%%%%%

Problem solving is rarely a perfectly linear process.

Nevertheless, a useful general framework is:

\begin{enumerate}
    \item understand the problem;
    \item organize the information;
    \item choose a strategy;
    \item carry out the strategy;
    \item check the result;
    \item communicate the solution.
\end{enumerate}

These stages may repeat.

A failed strategy often provides information that suggests a better one.

%%%%%%%%%%%%%%%%%%%%%%%%%%%%%%%%%%%%%%%%%%%%%%%%%%%%%%%%%%%%
\subsection{Step 1: Understand the Problem}
\label{subsec:understand-problem}
%%%%%%%%%%%%%%%%%%%%%%%%%%%%%%%%%%%%%%%%%%%%%%%%%%%%%%%%%%%%

Before calculating, determine what the problem actually asks.

Identify:

\begin{itemize}
    \item the known quantities;
    \item the unknown quantities;
    \item the assumptions;
    \item the domain;
    \item any restrictions;
    \item the desired conclusion.
\end{itemize}

A surprising number of errors begin with solving a different problem from the
one that was asked.

%%%%%%%%%%%%%%%%%%%%%%%%%%%%%%%%%%%%%%%%%%%%%%%%%%%%%%%%%%%%
\subsection{Identify the Goal}
\label{subsec:identify-goal}
%%%%%%%%%%%%%%%%%%%%%%%%%%%%%%%%%%%%%%%%%%%%%%%%%%%%%%%%%%%%

Different goals require different forms of answers.

For example:

\begin{itemize}
    \item ``Solve'' asks for values satisfying a condition.
    \item ``Prove'' asks for a logically complete argument.
    \item ``Show that'' generally asks for proof or derivation.
    \item ``Find'' asks for a specific mathematical object or value.
    \item ``Determine whether'' requires a decision together with justification.
    \item ``Give a counterexample'' asks for one example disproving a universal claim.
\end{itemize}

Understanding the verb in the problem is part of understanding the problem.

%%%%%%%%%%%%%%%%%%%%%%%%%%%%%%%%%%%%%%%%%%%%%%%%%%%%%%%%%%%%
\subsection{Step 2: Organize the Information}
\label{subsec:organize-information}
%%%%%%%%%%%%%%%%%%%%%%%%%%%%%%%%%%%%%%%%%%%%%%%%%%%%%%%%%%%%

Rewrite the given information in a useful form.

A verbal statement may become an equation.

A geometric problem may become a diagram.

A collection of numerical observations may become a table.

A logical condition may become an implication.

A complicated expression may be factored or rewritten.

Organization often reveals structure that was hidden in the original wording.

%%%%%%%%%%%%%%%%%%%%%%%%%%%%%%%%%%%%%%%%%%%%%%%%%%%%%%%%%%%%
\subsection{Step 3: Choose a Strategy}
\label{subsec:choose-strategy}
%%%%%%%%%%%%%%%%%%%%%%%%%%%%%%%%%%%%%%%%%%%%%%%%%%%%%%%%%%%%

Possible strategies include:

\begin{itemize}
    \item applying a definition;
    \item drawing a diagram;
    \item making a table;
    \item introducing variables;
    \item writing an equation;
    \item simplifying an expression;
    \item factoring;
    \item working backward;
    \item considering special cases;
    \item searching for a pattern;
    \item breaking the problem into smaller parts;
    \item using a known theorem;
    \item using contradiction or contraposition;
    \item constructing an example;
    \item looking for a counterexample.
\end{itemize}

Choosing a strategy is often the most creative part of mathematics.

%%%%%%%%%%%%%%%%%%%%%%%%%%%%%%%%%%%%%%%%%%%%%%%%%%%%%%%%%%%%
\subsection{Step 4: Carry Out the Strategy}
\label{subsec:carry-out-strategy}
%%%%%%%%%%%%%%%%%%%%%%%%%%%%%%%%%%%%%%%%%%%%%%%%%%%%%%%%%%%%

Once a strategy is selected, execute it carefully.

Keep track of:

\begin{itemize}
    \item algebraic restrictions;
    \item domains;
    \item signs;
    \item units;
    \item logical assumptions;
    \item whether transformations are reversible.
\end{itemize}

A correct strategy can still fail through careless execution.

%%%%%%%%%%%%%%%%%%%%%%%%%%%%%%%%%%%%%%%%%%%%%%%%%%%%%%%%%%%%
\subsection{Step 5: Check the Result}
\label{subsec:check-result}
%%%%%%%%%%%%%%%%%%%%%%%%%%%%%%%%%%%%%%%%%%%%%%%%%%%%%%%%%%%%

A solution should be checked whenever possible.

Useful checks include:

\begin{itemize}
    \item substitution into the original equation;
    \item checking units;
    \item estimating magnitude;
    \item checking signs;
    \item checking the domain;
    \item considering special cases;
    \item testing boundary values;
    \item comparing with known behavior.
\end{itemize}

Checking does not replace proof, but it is an important tool for detecting
errors.

%%%%%%%%%%%%%%%%%%%%%%%%%%%%%%%%%%%%%%%%%%%%%%%%%%%%%%%%%%%%
\subsection{Step 6: Communicate the Solution}
\label{subsec:communicate-solution}
%%%%%%%%%%%%%%%%%%%%%%%%%%%%%%%%%%%%%%%%%%%%%%%%%%%%%%%%%%%%

A finished solution should usually be more organized than the scratch work
used to discover it.

The final presentation should contain:

\begin{itemize}
    \item the essential reasoning;
    \item enough intermediate work to verify the result;
    \item definitions or theorems when needed;
    \item a clearly identified conclusion.
\end{itemize}

Exploration may be messy.

Presentation should not be.

%%%%%%%%%%%%%%%%%%%%%%%%%%%%%%%%%%%%%%%%%%%%%%%%%%%%%%%%%%%%
\subsection{Worked Example: A Complete Problem-Solving Process}
\label{subsec:worked-complete-problem-solving}
%%%%%%%%%%%%%%%%%%%%%%%%%%%%%%%%%%%%%%%%%%%%%%%%%%%%%%%%%%%%

\begin{workedexamplebox}{Solve and Check a Rational Equation}

Solve

\[ \frac{2}{x-1}=3. \]

\begin{tabularx}{\textwidth}{@{}>{\raggedright\arraybackslash}p{0.42\textwidth}X@{}}
Identify the domain restriction. & \(x\neq1\)\\
Multiply by \(x-1\). & \(2=3(x-1)\)\\
Distribute. & \(2=3x-3\)\\
Add \(3\). & \(5=3x\)\\
Divide by \(3\). & \(x=\frac53\)\\
Check the restriction. & \(\frac53\neq1\)\\
Substitute into the original equation. & \(\frac{2}{5/3-1}=3\)
\end{tabularx}

\begin{answerbox}
\[ x=\frac53 \]
\end{answerbox}
\end{workedexamplebox}

%%%%%%%%%%%%%%%%%%%%%%%%%%%%%%%%%%%%%%%%%%%%%%%%%%%%%%%%%%%%
\subsection{Work Forward}
\label{subsec:work-forward}
%%%%%%%%%%%%%%%%%%%%%%%%%%%%%%%%%%%%%%%%%%%%%%%%%%%%%%%%%%%%

Working forward means beginning with known information and determining what
follows.

Suppose

\[ x>3. \]

Then one may deduce

\[ x+2>5. \]

If \(x\) is known to be positive, one may also deduce

\[ x^2>0. \]

Forward reasoning asks:

\begin{center}
``What can I conclude from what I already know?''
\end{center}

%%%%%%%%%%%%%%%%%%%%%%%%%%%%%%%%%%%%%%%%%%%%%%%%%%%%%%%%%%%%
\subsection{Work Backward}
\label{subsec:work-backward}
%%%%%%%%%%%%%%%%%%%%%%%%%%%%%%%%%%%%%%%%%%%%%%%%%%%%%%%%%%%%

Working backward begins with the desired conclusion and asks what would be
sufficient to establish it.

Suppose the goal is

\[ (x-2)^2=0. \]

One may ask:

\begin{center}
What condition would make this true?
\end{center}

Clearly,

\[ x-2=0 \]

would be sufficient, which suggests

\[ x=2. \]

Backward reasoning is especially useful during problem discovery.

%%%%%%%%%%%%%%%%%%%%%%%%%%%%%%%%%%%%%%%%%%%%%%%%%%%%%%%%%%%%
\subsection{Meet in the Middle}
\label{subsec:meet-in-middle}
%%%%%%%%%%%%%%%%%%%%%%%%%%%%%%%%%%%%%%%%%%%%%%%%%%%%%%%%%%%%

Many difficult problems are solved by combining forward and backward
reasoning.

From the hypotheses, derive useful consequences.

From the goal, determine what intermediate result would be sufficient.

Then search for a connection between the two.

This strategy is common in proofs, geometry, algebra, and analysis.

%%%%%%%%%%%%%%%%%%%%%%%%%%%%%%%%%%%%%%%%%%%%%%%%%%%%%%%%%%%%
\subsection{Break a Problem into Smaller Parts}
\label{subsec:break-problem-smaller}
%%%%%%%%%%%%%%%%%%%%%%%%%%%%%%%%%%%%%%%%%%%%%%%%%%%%%%%%%%%%

A complicated problem often becomes manageable when decomposed.

For example, to analyze

\[ f(x)=\frac{x^2-4}{x^2-x-6}, \]

one might separately:

\begin{enumerate}
    \item factor the numerator;
    \item factor the denominator;
    \item determine domain restrictions;
    \item simplify common factors;
    \item interpret the remaining expression.
\end{enumerate}

Each smaller task is easier than handling the entire expression at once.

%%%%%%%%%%%%%%%%%%%%%%%%%%%%%%%%%%%%%%%%%%%%%%%%%%%%%%%%%%%%
\subsection{Draw a Diagram}
\label{subsec:draw-diagram}
%%%%%%%%%%%%%%%%%%%%%%%%%%%%%%%%%%%%%%%%%%%%%%%%%%%%%%%%%%%%

A diagram can reveal relationships that are difficult to see symbolically.

Diagrams are particularly useful in:

\begin{itemize}
    \item geometry;
    \item trigonometry;
    \item vector analysis;
    \item graph theory;
    \item probability;
    \item optimization;
    \item multivariable calculus.
\end{itemize}

A diagram should support reasoning rather than replace it.

A picture may suggest that something is true without proving it.

%%%%%%%%%%%%%%%%%%%%%%%%%%%%%%%%%%%%%%%%%%%%%%%%%%%%%%%%%%%%
\subsection{Make a Table}
\label{subsec:make-table}
%%%%%%%%%%%%%%%%%%%%%%%%%%%%%%%%%%%%%%%%%%%%%%%%%%%%%%%%%%%%

Tables organize repeated calculations and can reveal patterns.

Suppose

\[ a_n=n^2-n. \]

Evaluating several values gives

\[
\begin{array}{c|ccccc}
n&1&2&3&4&5\\ \hline
a_n&0&2&6&12&20
\end{array}
\]

The values suggest

\[ a_n=n(n-1), \]

which is always even because one of two consecutive integers is even.

The table helps discover the pattern; a proof establishes it.

%%%%%%%%%%%%%%%%%%%%%%%%%%%%%%%%%%%%%%%%%%%%%%%%%%%%%%%%%%%%
\subsection{Search for Patterns}
\label{subsec:search-patterns}
%%%%%%%%%%%%%%%%%%%%%%%%%%%%%%%%%%%%%%%%%%%%%%%%%%%%%%%%%%%%

Patterns often suggest conjectures.

For example,

\[ 1=1^2, \]

\[ 1+3=4=2^2, \]

\[ 1+3+5=9=3^2, \]

\[ 1+3+5+7=16=4^2. \]

This suggests

\[ 1+3+\cdots+(2n-1)=n^2. \]

The observed pattern motivates the statement.

Mathematical induction can then prove it.

%%%%%%%%%%%%%%%%%%%%%%%%%%%%%%%%%%%%%%%%%%%%%%%%%%%%%%%%%%%%
\subsection{Conjectures}
\label{subsec:problem-solving-conjectures}
%%%%%%%%%%%%%%%%%%%%%%%%%%%%%%%%%%%%%%%%%%%%%%%%%%%%%%%%%%%%

A \emph{conjecture} is a mathematical statement believed to be true based on
evidence or reasoning but not yet established by proof.

Experimentation is valuable for producing conjectures.

However,

\begin{importantbox}
A pattern suggests a theorem; it does not prove a theorem.
\end{importantbox}

%%%%%%%%%%%%%%%%%%%%%%%%%%%%%%%%%%%%%%%%%%%%%%%%%%%%%%%%%%%%
\subsection{Test Special Cases}
\label{subsec:test-special-cases}
%%%%%%%%%%%%%%%%%%%%%%%%%%%%%%%%%%%%%%%%%%%%%%%%%%%%%%%%%%%%

Special cases can clarify a general problem.

If a statement involves a parameter \(n\), test values such as

\[ n=1,\qquad n=2,\qquad n=3. \]

If it involves a real variable, useful cases may include

\[ x=0,\qquad x=1,\qquad x=-1. \]

Special cases can:

\begin{itemize}
    \item reveal errors;
    \item suggest patterns;
    \item identify missing assumptions;
    \item produce counterexamples.
\end{itemize}

%%%%%%%%%%%%%%%%%%%%%%%%%%%%%%%%%%%%%%%%%%%%%%%%%%%%%%%%%%%%
\subsection{Look for Counterexamples}
\label{subsec:look-counterexamples}
%%%%%%%%%%%%%%%%%%%%%%%%%%%%%%%%%%%%%%%%%%%%%%%%%%%%%%%%%%%%

Before attempting to prove a statement, ask whether it is actually true.

Consider the claim

\[ \sqrt{x^2}=x. \]

Testing

\[ x=-3 \]

gives

\[ \sqrt{(-3)^2}=3\neq-3. \]

Thus the claim is false.

The correct identity is

\[ \sqrt{x^2}=|x|. \]

A small test can prevent a long attempt to prove a false statement.

%%%%%%%%%%%%%%%%%%%%%%%%%%%%%%%%%%%%%%%%%%%%%%%%%%%%%%%%%%%%
\subsection{Estimate Before Calculating}
\label{subsec:estimate-before-calculating}
%%%%%%%%%%%%%%%%%%%%%%%%%%%%%%%%%%%%%%%%%%%%%%%%%%%%%%%%%%%%

Estimation gives a rough expectation for the answer.

For example,

\[ 19.8(5.1) \]

should be close to

\[ 20(5)=100. \]

If a calculator returns

\[ 10.098, \]

the result should immediately appear suspicious.

Estimation provides a simple error-detection mechanism.

%%%%%%%%%%%%%%%%%%%%%%%%%%%%%%%%%%%%%%%%%%%%%%%%%%%%%%%%%%%%
\subsection{Exact and Approximate Answers}
\label{subsec:exact-approximate-writing}
%%%%%%%%%%%%%%%%%%%%%%%%%%%%%%%%%%%%%%%%%%%%%%%%%%%%%%%%%%%%

Whenever practical, preserve exact values until approximation is needed.

For example,

\[ \sqrt2 \]

is exact, whereas

\[ 1.414 \]

is approximate.

The symbol

\[ \approx \]

is read

\begin{center}
``is approximately equal to.''
\end{center}

Thus one may write

\[ \sqrt2\approx1.414. \]

Using \(=\) instead of \(\approx\) would incorrectly claim exact equality.

%%%%%%%%%%%%%%%%%%%%%%%%%%%%%%%%%%%%%%%%%%%%%%%%%%%%%%%%%%%%
\subsection{Check Units}
\label{subsec:check-units}
%%%%%%%%%%%%%%%%%%%%%%%%%%%%%%%%%%%%%%%%%%%%%%%%%%%%%%%%%%%%

Units provide mathematical information.

Suppose distance is measured in meters and time in seconds.

Then

\[ \frac{\text{distance}}{\text{time}} \]

has units

\[ \frac{\mathrm m}{\mathrm s}. \]

A proposed speed with units

\[ \mathrm m^2 \]

would therefore signal an error.

This type of reasoning is called \emph{dimensional analysis} in many applied
contexts.

%%%%%%%%%%%%%%%%%%%%%%%%%%%%%%%%%%%%%%%%%%%%%%%%%%%%%%%%%%%%
\subsection{Unit Consistency}
\label{subsec:unit-consistency}
%%%%%%%%%%%%%%%%%%%%%%%%%%%%%%%%%%%%%%%%%%%%%%%%%%%%%%%%%%%%

Quantities being added or subtracted should have compatible units.

For example,

\[ 3\text{ meters}+5\text{ meters}=8\text{ meters} \]

is meaningful.

But

\[ 3\text{ meters}+5\text{ seconds} \]

does not represent an ordinary physical sum.

Unit consistency becomes especially important in applied mathematics,
physics, engineering, and mathematical modeling.

%%%%%%%%%%%%%%%%%%%%%%%%%%%%%%%%%%%%%%%%%%%%%%%%%%%%%%%%%%%%
\subsection{Check the Domain}
\label{subsec:check-domain}
%%%%%%%%%%%%%%%%%%%%%%%%%%%%%%%%%%%%%%%%%%%%%%%%%%%%%%%%%%%%

Algebraic manipulation may produce values that are not allowed in the
original problem.

Consider

\[ \frac{x+1}{x-2}=0. \]

The denominator requires

\[ x\neq2. \]

The equation is zero when the numerator is zero:

\[ x=-1. \]

The domain restriction must remain part of the reasoning.

%%%%%%%%%%%%%%%%%%%%%%%%%%%%%%%%%%%%%%%%%%%%%%%%%%%%%%%%%%%%
\subsection{Extraneous Solutions}
\label{subsec:extraneous-solutions}
%%%%%%%%%%%%%%%%%%%%%%%%%%%%%%%%%%%%%%%%%%%%%%%%%%%%%%%%%%%%

Some operations can introduce candidate solutions that do not satisfy the
original equation.

For example, squaring both sides is not always reversible.

If

\[ x=-2, \]

then

\[ x^2=4. \]

But starting from

\[ x^2=4 \]

does not allow one to conclude only \(x=-2\); both signs are possible.

In more complicated equations, squaring can produce extraneous solutions.

Therefore candidate solutions should often be checked in the original
equation.

%%%%%%%%%%%%%%%%%%%%%%%%%%%%%%%%%%%%%%%%%%%%%%%%%%%%%%%%%%%%
\subsection{Reversible and Irreversible Steps}
\label{subsec:reversible-steps}
%%%%%%%%%%%%%%%%%%%%%%%%%%%%%%%%%%%%%%%%%%%%%%%%%%%%%%%%%%%%

An equivalence transformation preserves exactly the same solution set.

For example,

\[ x+3=7\Longleftrightarrow x=4. \]

However, an implication may preserve only one direction.

For example,

\[ x=2\Longrightarrow x^2=4, \]

but

\[ x^2=4\not\Longrightarrow x=2. \]

The distinction between implication and equivalence is important when solving
equations and writing proofs.

%%%%%%%%%%%%%%%%%%%%%%%%%%%%%%%%%%%%%%%%%%%%%%%%%%%%%%%%%%%%
\subsection{Reading Advanced Notation}
\label{subsec:reading-advanced-notation}
%%%%%%%%%%%%%%%%%%%%%%%%%%%%%%%%%%%%%%%%%%%%%%%%%%%%%%%%%%%%

As mathematics becomes more advanced, notation becomes increasingly compact.

The goal should remain to understand the mathematical sentence represented by
the symbols.

For example, later in calculus one encounters

\[ \lim_{x\to a}f(x)=L. \]

The notation

\[ \lim \]

is read ``limit.''

In this context, the arrow

\[ x\to a \]

is read

\begin{center}
``\(x\) approaches \(a\).''
\end{center}

Therefore,

\[ \lim_{x\to a}f(x) \]

is read

\begin{center}
``the limit of \(f(x)\) as \(x\) approaches \(a\).''
\end{center}

The complete statement

\[ \lim_{x\to a}f(x)=L \]

is read

\begin{center}
``the limit of \(f(x)\) as \(x\) approaches \(a\) equals \(L\).''
\end{center}

The rigorous meaning of limits is developed later in Analysis.

%%%%%%%%%%%%%%%%%%%%%%%%%%%%%%%%%%%%%%%%%%%%%%%%%%%%%%%%%%%%
\subsection{Context Determines How Symbols Are Read}
\label{subsec:context-symbol-reading}
%%%%%%%%%%%%%%%%%%%%%%%%%%%%%%%%%%%%%%%%%%%%%%%%%%%%%%%%%%%%

The same symbol can have different readings in different mathematical
contexts.

For example,

\[ P\to Q \]

is read

\begin{center}
``\(P\) implies \(Q\).''
\end{center}

For a function,

\[ f:A\to B \]

is read

\begin{center}
``\(f\) maps \(A\) to \(B\).''
\end{center}

In calculus,

\[ x\to a \]

is read

\begin{center}
``\(x\) approaches \(a\).''
\end{center}

The shape of a symbol alone does not determine its interpretation.

%%%%%%%%%%%%%%%%%%%%%%%%%%%%%%%%%%%%%%%%%%%%%%%%%%%%%%%%%%%%
\subsection{Writing Proofs Clearly}
\label{subsec:writing-proofs-clearly}
%%%%%%%%%%%%%%%%%%%%%%%%%%%%%%%%%%%%%%%%%%%%%%%%%%%%%%%%%%%%

A proof should communicate its logical structure.

For example:

\begin{center}
Let \(n\) be an even integer. By definition, \(n=2k\) for some \(k\in\Z\).
Then
\[ n^2=(2k)^2=4k^2=2(2k^2). \]
Since \(2k^2\in\Z\), \(n^2\) is even.
\end{center}

The proof contains:

\begin{itemize}
    \item an arbitrary object;
    \item the relevant definition;
    \item the calculation;
    \item the final connection back to the definition.
\end{itemize}

%%%%%%%%%%%%%%%%%%%%%%%%%%%%%%%%%%%%%%%%%%%%%%%%%%%%%%%%%%%%
\subsection{Do Not Hide the Conclusion}
\label{subsec:do-not-hide-conclusion}
%%%%%%%%%%%%%%%%%%%%%%%%%%%%%%%%%%%%%%%%%%%%%%%%%%%%%%%%%%%%

A proof should clearly indicate when the desired result has been established.

Phrases such as

\begin{center}
``Therefore \(n^2\) is even''
\end{center}

or

\begin{center}
``Hence \(A=B\)''
\end{center}

help the reader connect the final step with the original claim.

%%%%%%%%%%%%%%%%%%%%%%%%%%%%%%%%%%%%%%%%%%%%%%%%%%%%%%%%%%%%
\subsection{Writing Worked Solutions}
\label{subsec:writing-worked-solutions}
%%%%%%%%%%%%%%%%%%%%%%%%%%%%%%%%%%%%%%%%%%%%%%%%%%%%%%%%%%%%

A good worked solution should answer three questions:

\begin{enumerate}
    \item What are we trying to find?
    \item Why is each important step valid?
    \item What is the final answer?
\end{enumerate}

The amount of explanation depends on the problem and intended audience.

Elementary algebra may require more intermediate steps than an advanced
analysis calculation.

%%%%%%%%%%%%%%%%%%%%%%%%%%%%%%%%%%%%%%%%%%%%%%%%%%%%%%%%%%%%
\subsection{Worked Example: Improving a Solution}
\label{subsec:worked-improving-solution}
%%%%%%%%%%%%%%%%%%%%%%%%%%%%%%%%%%%%%%%%%%%%%%%%%%%%%%%%%%%%

\begin{workedexamplebox}{Write a Clear Solution}

Solve

\[ 4(x-3)+5=17. \]

\begin{tabularx}{\textwidth}{@{}>{\raggedright\arraybackslash}p{0.42\textwidth}X@{}}
Distribute \(4\). & \(4x-12+5=17\)\\
Combine like terms. & \(4x-7=17\)\\
Add \(7\) to both sides. & \(4x=24\)\\
Divide by \(4\). & \(x=6\)\\
Check. & \(4(6-3)+5=17\)
\end{tabularx}

\begin{answerbox}
\[ x=6 \]
\end{answerbox}
\end{workedexamplebox}

%%%%%%%%%%%%%%%%%%%%%%%%%%%%%%%%%%%%%%%%%%%%%%%%%%%%%%%%%%%%
\subsection{Exploration versus Presentation}
\label{subsec:exploration-vs-presentation}
%%%%%%%%%%%%%%%%%%%%%%%%%%%%%%%%%%%%%%%%%%%%%%%%%%%%%%%%%%%%

The process of discovering mathematics is often messy.

Scratch work may include:

\begin{itemize}
    \item failed approaches;
    \item numerical experiments;
    \item diagrams;
    \item guesses;
    \item incomplete calculations;
    \item rewritten expressions.
\end{itemize}

These are valuable parts of problem solving.

However, the final solution should usually present the successful argument in
a logical order.

\begin{conceptbox}
Discovery and presentation serve different purposes.

Exploration helps the mathematician find the solution.

Presentation helps the reader understand the solution.
\end{conceptbox}

%%%%%%%%%%%%%%%%%%%%%%%%%%%%%%%%%%%%%%%%%%%%%%%%%%%%%%%%%%%%
\subsection{When a Strategy Fails}
\label{subsec:when-strategy-fails}
%%%%%%%%%%%%%%%%%%%%%%%%%%%%%%%%%%%%%%%%%%%%%%%%%%%%%%%%%%%%

A failed strategy is not necessarily wasted effort.

It may reveal:

\begin{itemize}
    \item which information is insufficient;
    \item which assumptions matter;
    \item which definitions should be used;
    \item which cases behave differently;
    \item whether the statement may be false.
\end{itemize}

Problem solving is often iterative rather than immediate.

%%%%%%%%%%%%%%%%%%%%%%%%%%%%%%%%%%%%%%%%%%%%%%%%%%%%%%%%%%%%
\subsection{Use Technology as a Tool}
\label{subsec:technology-problem-solving}
%%%%%%%%%%%%%%%%%%%%%%%%%%%%%%%%%%%%%%%%%%%%%%%%%%%%%%%%%%%%

Calculators, computer algebra systems, numerical software, graphing programs,
and programming languages can assist mathematical exploration.

They can help:

\begin{itemize}
    \item calculate examples;
    \item visualize functions;
    \item test conjectures;
    \item approximate solutions;
    \item simulate mathematical models;
    \item manipulate large datasets.
\end{itemize}

However, computational output should still be interpreted mathematically.

%%%%%%%%%%%%%%%%%%%%%%%%%%%%%%%%%%%%%%%%%%%%%%%%%%%%%%%%%%%%
\subsection{Computation Is Not Automatically Proof}
\label{subsec:computation-not-proof}
%%%%%%%%%%%%%%%%%%%%%%%%%%%%%%%%%%%%%%%%%%%%%%%%%%%%%%%%%%%%

Suppose a computer checks a statement for the first million positive integers.

This is strong evidence.

It does not generally prove the statement for every positive integer.

Similarly, a graph may suggest that two functions intersect once, but the
display alone may not establish uniqueness.

Technology supports reasoning; it does not automatically replace it.

%%%%%%%%%%%%%%%%%%%%%%%%%%%%%%%%%%%%%%%%%%%%%%%%%%%%%%%%%%%%
\subsection{Common Mathematical Writing Problems}
\label{subsec:common-writing-problems}
%%%%%%%%%%%%%%%%%%%%%%%%%%%%%%%%%%%%%%%%%%%%%%%%%%%%%%%%%%%%

\subsubsection{Undefined Variables}

Poor:

\[ x^2+y^2=r^2. \]

Better:

\begin{center}
Let \(x\) and \(y\) denote the legs of a right triangle and let \(r\) denote
the length of its hypotenuse. Then
\[ x^2+y^2=r^2. \]
\end{center}

%%%%%%%%%%%%%%%%%%%%%%%%%%%%%%%%%%%%%%%%%%%%%%%%%%%%%%%%%%%%
\subsubsection{Ambiguous Pronouns}
\label{subsec:ambiguous-pronouns}
%%%%%%%%%%%%%%%%%%%%%%%%%%%%%%%%%%%%%%%%%%%%%%%%%%%%%%%%%%%%

Statements such as

\begin{center}
``This shows it is true''
\end{center}

may be unclear when several objects have just been discussed.

Name the object or conclusion explicitly.

%%%%%%%%%%%%%%%%%%%%%%%%%%%%%%%%%%%%%%%%%%%%%%%%%%%%%%%%%%%%
\subsubsection{Unexplained Equations}
\label{subsec:unexplained-equations}
%%%%%%%%%%%%%%%%%%%%%%%%%%%%%%%%%%%%%%%%%%%%%%%%%%%%%%%%%%%%

A long chain of equations without explanation may leave the reader unable to
identify the purpose of the calculation.

Use brief prose to identify important steps.

%%%%%%%%%%%%%%%%%%%%%%%%%%%%%%%%%%%%%%%%%%%%%%%%%%%%%%%%%%%%
\subsubsection{Too Much Explanation}
\label{subsec:too-much-explanation}
%%%%%%%%%%%%%%%%%%%%%%%%%%%%%%%%%%%%%%%%%%%%%%%%%%%%%%%%%%%%

Clarity does not require explaining every trivial arithmetic operation.

Writing

\begin{center}
``Now add \(1\) to \(2\), which produces \(3\)''
\end{center}

may be unnecessary in an advanced text.

The appropriate level of explanation depends on the mathematical difficulty
and intended audience.

%%%%%%%%%%%%%%%%%%%%%%%%%%%%%%%%%%%%%%%%%%%%%%%%%%%%%%%%%%%%
\subsubsection{Incorrect Use of Equality}
\label{subsec:incorrect-equality-writing}
%%%%%%%%%%%%%%%%%%%%%%%%%%%%%%%%%%%%%%%%%%%%%%%%%%%%%%%%%%%%

Do not use \(=\) to mean:

\begin{itemize}
    \item ``therefore'';
    \item ``approximately'';
    \item ``the next step is'';
    \item ``implies.''
\end{itemize}

Each of these ideas has its own mathematical or linguistic expression.

%%%%%%%%%%%%%%%%%%%%%%%%%%%%%%%%%%%%%%%%%%%%%%%%%%%%%%%%%%%%
\subsubsection{Ignoring Restrictions}
\label{subsec:ignoring-restrictions-writing}
%%%%%%%%%%%%%%%%%%%%%%%%%%%%%%%%%%%%%%%%%%%%%%%%%%%%%%%%%%%%

If

\[ f(x)=\frac1{x-3}, \]

then

\[ x\neq3 \]

is part of the mathematical information.

A solution that ignores the restriction is incomplete.

%%%%%%%%%%%%%%%%%%%%%%%%%%%%%%%%%%%%%%%%%%%%%%%%%%%%%%%%%%%%
\subsubsection{Premature Decimal Approximation}
\label{subsec:premature-rounding}
%%%%%%%%%%%%%%%%%%%%%%%%%%%%%%%%%%%%%%%%%%%%%%%%%%%%%%%%%%%%

Replacing exact values by rounded decimals too early can introduce accumulated
error.

For example, it is usually preferable to retain

\[ \frac{\sqrt2}{3} \]

during symbolic work and approximate only when a decimal answer is required.

%%%%%%%%%%%%%%%%%%%%%%%%%%%%%%%%%%%%%%%%%%%%%%%%%%%%%%%%%%%%
\subsection{A General Problem-Solving Checklist}
\label{subsec:problem-solving-checklist}
%%%%%%%%%%%%%%%%%%%%%%%%%%%%%%%%%%%%%%%%%%%%%%%%%%%%%%%%%%%%

Before considering a solution complete, ask:

\begin{enumerate}
    \item What exactly was asked?
    \item Did I identify all assumptions and restrictions?
    \item Are my variables clearly defined?
    \item Did I use the correct domain?
    \item Is each important step justified?
    \item Did I accidentally divide by zero or apply an invalid operation?
    \item Are the units consistent?
    \item Is the magnitude reasonable?
    \item Did I check possible extraneous solutions?
    \item Does the final answer actually answer the original question?
\end{enumerate}

%%%%%%%%%%%%%%%%%%%%%%%%%%%%%%%%%%%%%%%%%%%%%%%%%%%%%%%%%%%%
\subsection{Worked Example: Diagnose an Incorrect Solution}
\label{subsec:worked-diagnose-solution}
%%%%%%%%%%%%%%%%%%%%%%%%%%%%%%%%%%%%%%%%%%%%%%%%%%%%%%%%%%%%

\begin{workedexamplebox}{Find the Error}

A student solves

\[ \sqrt{x+6}=x \]

by squaring:

\[ x+6=x^2. \]

Rearranging gives

\[ x^2-x-6=0, \]

so

\[ (x-3)(x+2)=0. \]

The candidate solutions are

\[ x=3\qquad\text{and}\qquad x=-2. \]

\begin{tabularx}{\textwidth}{@{}>{\raggedright\arraybackslash}p{0.42\textwidth}X@{}}
Check \(x=3\). & \(\sqrt{9}=3\), so it works\\
Check \(x=-2\). & \(\sqrt4=2\neq-2\), so it fails\\
Squaring can introduce extra solutions. & \(x=-2\) is extraneous
\end{tabularx}

\begin{answerbox}
\[ x=3 \]
\end{answerbox}
\end{workedexamplebox}

%%%%%%%%%%%%%%%%%%%%%%%%%%%%%%%%%%%%%%%%%%%%%%%%%%%%%%%%%%%%
\subsection{Worked Example: Translate Before Solving}
\label{subsec:worked-translate-solving}
%%%%%%%%%%%%%%%%%%%%%%%%%%%%%%%%%%%%%%%%%%%%%%%%%%%%%%%%%%%%

\begin{workedexamplebox}{Translate a Verbal Problem}

A rectangle has length \(3\) meters greater than its width and perimeter
\(30\) meters. Find its dimensions.

Let

\[ w=\text{width}. \]

Then

\[ \ell=w+3, \]

where the symbol \(\ell\) is the lowercase Greek letter \emph{ell} used here
to represent length and is read ``ell.''

The perimeter equation is

\[ 2\ell+2w=30. \]

\begin{tabularx}{\textwidth}{@{}>{\raggedright\arraybackslash}p{0.42\textwidth}X@{}}
Substitute \(\ell=w+3\). & \(2(w+3)+2w=30\)\\
Distribute. & \(2w+6+2w=30\)\\
Combine terms. & \(4w+6=30\)\\
Solve for \(w\). & \(w=6\)\\
Find the length. & \(\ell=6+3=9\)
\end{tabularx}

\begin{answerbox}
The rectangle is \(6\) meters wide and \(9\) meters long.
\end{answerbox}
\end{workedexamplebox}

%%%%%%%%%%%%%%%%%%%%%%%%%%%%%%%%%%%%%%%%%%%%%%%%%%%%%%%%%%%%
\subsection{Exercises}
\label{subsec:mathematical-writing-exercises}
%%%%%%%%%%%%%%%%%%%%%%%%%%%%%%%%%%%%%%%%%%%%%%%%%%%%%%%%%%%%

\begin{exercise}[1]
Rewrite the following as a complete mathematical sentence:
\[ x>4,\qquad x^2>16. \]
\end{exercise}

\begin{exercise}[1]
Explain in words what
\[ x\in A\cap B \]
means.
\end{exercise}

\begin{exercise}[1]
Read aloud:
\[ f:A\to B. \]
\end{exercise}

\begin{exercise}[1]
Read aloud:
\[ \forall x\in\R,\;x^2\geq0. \]
\end{exercise}

\begin{exercise}[1]
Read aloud:
\[ \lim_{x\to a}f(x)=L. \]
\end{exercise}

\begin{exercise}[1]
Explain why
\[ \sqrt2=1.414 \]
is not mathematically exact. Rewrite it correctly.
\end{exercise}

\begin{exercise}[1]
Identify the domain restriction in
\[ f(x)=\frac{x+2}{x-7}. \]
\end{exercise}

\begin{exercise}[1]
Explain what is wrong with the chain
\[ x^2=9=x=3. \]
\end{exercise}

\begin{exercise}[2]
Solve and check:
\[ \frac{3}{x+2}=6. \]
\end{exercise}

\begin{exercise}[2]
Solve and check:
\[ \sqrt{x+5}=x-1. \]
\end{exercise}

\begin{exercise}[2]
A rectangle has length \(5\) units greater than its width and perimeter
\(42\) units. Define variables, write an equation, and determine the
dimensions.
\end{exercise}

\begin{exercise}[2]
A car travels \(180\) miles in \(3\) hours. Determine its average speed and
include the correct units.
\end{exercise}

\begin{exercise}[2]
Before calculating exactly, estimate
\[ 49.7(2.03). \]
Then calculate and compare.
\end{exercise}

\begin{exercise}[2]
Consider the claim
\[ \sqrt{a+b}=\sqrt a+\sqrt b. \]
Test special cases and determine whether the statement is always true.
\end{exercise}

\begin{exercise}[2]
Consider the sequence
\[ 2,6,12,20,30,\ldots. \]
Compute several related expressions and propose a formula for the \(n\)-th
term.
\end{exercise}

\begin{exercise}[2]
Explain why checking the first \(1000\) cases of a universal statement does
not generally constitute a proof.
\end{exercise}

\begin{exercise}[2]
Write a clear proof that the sum of two odd integers is even. Introduce every
variable and state the relevant definitions.
\end{exercise}

\begin{exercise}[2]
Rewrite the following argument more clearly:
\[ n=2k+1,\quad m=2j+1,\quad n+m=2k+2j+2=2(k+j+1),\quad\text{even}. \]
\end{exercise}

\begin{exercise}[2]
Explain why the transformation
\[ x=3\Longrightarrow x^2=9 \]
is valid but not reversible without additional possibilities.
\end{exercise}

\begin{exercise}[2]
A calculation produces a length of \(-4\) meters. Explain why the mathematical
context should cause the result to be reconsidered.
\end{exercise}

\begin{exercise}[2]
A speed calculation produces units of
\[ \mathrm{m}\cdot\mathrm{s}. \]
Explain why dimensional analysis suggests an error.
\end{exercise}

\begin{exercise}[3]
Develop a conjecture from the pattern
\[ 1=1,\qquad 1+2=3,\qquad 1+2+3=6,\qquad 1+2+3+4=10. \]
Then state what additional work would be needed to turn the conjecture into a
theorem.
\end{exercise}

\begin{exercise}[3]
Suppose a computer verifies a number-theoretic claim for every integer from
\(1\) through \(10^9\). Explain what this establishes and what it does not
establish.
\end{exercise}

\begin{exercise}[3]
Choose a theorem from \cref{sec:methods-of-proof}. Identify its hypotheses,
conclusion, definitions used, and overall proof strategy.
\end{exercise}

%%%%%%%%%%%%%%%%%%%%%%%%%%%%%%%%%%%%%%%%%%%%%%%%%%%%%%%%%%%%
\subsection{Chapter 0 Connections}
\label{subsec:foundations-final-connections}
%%%%%%%%%%%%%%%%%%%%%%%%%%%%%%%%%%%%%%%%%%%%%%%%%%%%%%%%%%%%

\chapterconnection{
This section brings together the major ideas of Chapter 0. Mathematical
notation provides the language; arithmetic and number systems provide basic
objects and operations; set theory organizes collections; relations and
functions describe mathematical structure; logic formalizes statements;
proof methods establish truth; and cardinality extends the idea of size to
infinite sets. These tools form the common foundation for the branches of
mathematics developed throughout the remainder of this text.
}

%%%%%%%%%%%%%%%%%%%%%%%%%%%%%%%%%%%%%%%%%%%%%%%%%%%%%%%%%%%%
\subsection{Looking Ahead}
\label{subsec:foundations-looking-ahead}
%%%%%%%%%%%%%%%%%%%%%%%%%%%%%%%%%%%%%%%%%%%%%%%%%%%%%%%%%%%%

The remainder of this text develops increasingly specialized mathematical
structures, but the habits established in Chapter 0 remain essential.

When encountering a new concept:

\begin{enumerate}
    \item identify the objects being studied;
    \item learn the definitions precisely;
    \item understand the notation and how it is read;
    \item study examples and counterexamples;
    \item identify the central theorems;
    \item understand why those theorems are true;
    \item practice applying them to new problems;
    \item connect the new ideas to previously established mathematics.
\end{enumerate}

These habits turn mathematical study from memorizing isolated formulas into
understanding connected structures.

%%%%%%%%%%%%%%%%%%%%%%%%%%%%%%%%%%%%%%%%%%%%%%%%%%%%%%%%%%%%
% SECTION SOURCE TRACKING
%%%%%%%%%%%%%%%%%%%%%%%%%%%%%%%%%%%%%%%%%%%%%%%%%%%%%%%%%%%%

%------------------------------------------------------------
% 0.10 Mathematical Writing and Problem Solving
%------------------------------------------------------------
%
% Sources:
%   - Richard Hammack, Book of Proof, Third Edition.
%     Key: hammack2018bookofproof
%
% Used for:
%   - mathematical reading and writing conventions;
%   - introducing variables and mathematical objects;
%   - use of definitions in mathematical reasoning;
%   - organization and communication of proofs;
%   - distinction between examples, evidence, and proof;
%   - general proof-writing practices.
%
% Notes:
%   - General problem-solving strategies, worked examples, notation-reading
%     guidance, dimensional checks, and computational discussion are
%     independently developed for this text.
%   - Mathematical writing conventions continue throughout all later chapters.
%
\nocite{hammack2018bookofproof}